% --- Archon physics formalization source begin ---
% archon:physics
% archon:covers PhyXMiniProblems/problem_phyx_mini_0186.lean
% archon:source-report reports/phyx_mini/problem_phyx_mini_0186.source.json

\chapter{Physics problem phyx\_mini\_0186}
\label{ch:PhyXMiniProblems_problem_phyx_mini_0186}

\paragraph{Problem source.}
\#\# Physical scenario

The three identical loudspeakers in figure play a \$170 \textbackslash{}, \textbackslash{}text\{Hz\}\$ tone in a room where the speed of sound is \$340 \textbackslash{}, \textbackslash{}text\{m/s\}\$. You are standing \$4.0 \textbackslash{}, \textbackslash{}text\{m\}\$ in front of the middle speaker. At this point, the amplitude of the wave from each speaker is \$a\$.

\#\# Question

How far must speaker 2 be moved to the left to produce a maximum amplitude at the point where you are standing?

\#\# Answer choices

- A: A: 0.50 m

- B: B: 2.00 m

- C: C: 0.25 m

- D: D: 1.00 m

\#\# Figure caption (auxiliary; use the image as primary evidence)

The image depicts three audio speakers arranged vertically in a straight line. Each speaker is numbered (1, 2, and 3) from top to bottom. They emit sound waves, illustrated by blue curved lines emanating from the speakers.

The vertical distance between each speaker is labeled as 3.0 meters. To the right of the speakers, there is a black dot, which may represent a point of interest or measurement location. A horizontal arrow below the speakers indicates a distance of 4.0 meters from the black dot to the vertical line on which the speakers are positioned.

The arrangement suggests a setup possibly for sound wave interference or acoustics study, focusing on the distances between speakers and a point in the field.

\#\# Dataset metadata

Resonance and Harmonics; Physical Model Grounding Reasoning, Spatial Relation Reasoning

\paragraph{Recorded answer/context.}
D: D: 1.00 m

\paragraph{Figure/image path.}
/root/proposal\_for\_physic/science-mango/phyx\_mini\_run/phyx\_data/test\_image/186.png

\paragraph{Formalization target.}
create a compiling Lean file with sorry bodies at `PhyXMiniProblems/problem\_phyx\_mini\_0186.lean`. The Lean declarations must preserve the physical quantities, dimensions or dimensional roles, figure labels, governing-law hypotheses, and final relation expressed by this problem.
Use Mathlib/Physlib names found through LeanExplore where available. If a physics API is missing, introduce faithful local abstractions rather than scalar placeholder aliases.

\begin{theorem}[Physics formalization target]
\label{thm:physics:phyx_mini_0186:target}
\lean{PhyXMiniProblems.ProblemPhyXMini0186.problem_phyx_mini_0186}
\uses{def:physics:phyx-mini-0186:phyxminiproblems-problemphyxmini0186-lengthofmeters, def:physics:phyx-mini-0186:phyxminiproblems-problemphyxmini0186-threespeakerinterferencesetup, def:physics:phyx-mini-0186:phyxminiproblems-problemphyxmini0186-matchesproblemreadouts, def:physics:phyx-mini-0186:phyxminiproblems-problemphyxmini0186-hasdepictedlayout, def:physics:phyx-mini-0186:phyxminiproblems-problemphyxmini0186-emitscoherentlywithcommonamplitude, def:physics:phyx-mini-0186:phyxminiproblems-problemphyxmini0186-hasphysicalparameters, def:physics:phyx-mini-0186:phyxminiproblems-problemphyxmini0186-satisfiessoundwavelaw, def:physics:phyx-mini-0186:phyxminiproblems-problemphyxmini0186-isleastpositivemaximumdisplacement, def:physics:phyx-mini-0186:phyxminiproblems-problemphyxmini0186-iscorrectanswer, def:physics:phyx-mini-0186:phyxminiproblems-problemphyxmini0186-recordedanswerchoice}
The assigned autoformalize agent should translate this physics problem into Lean declarations in the covered file, with theorem and lemma proof bodies written as `by sorry`.
\end{theorem}
\begin{proof}
This is an autoformalization task, not a proof task. Produce statements that can later be proved by the physics prover without weakening the physical model.
\end{proof}

% --- Archon declaration topology begin ---
\section{Declaration topology}
\begin{definition}[Acoustic Length]
  \label{def:physics:phyx-mini-0186:phyxminiproblems-problemphyxmini0186-acousticlength}
  \lean{PhyXMiniProblems.ProblemPhyXMini0186.AcousticLength}
  A physical length, independent of the unit system used to read it
\end{definition}
\begin{proof}
  This is the declared model definition of Acoustic Length.
\end{proof}
\begin{definition}[Acoustic Frequency]
  \label{def:physics:phyx-mini-0186:phyxminiproblems-problemphyxmini0186-acousticfrequency}
  \lean{PhyXMiniProblems.ProblemPhyXMini0186.AcousticFrequency}
  A physical acoustic frequency, carrying the inverse-time dimension
\end{definition}
\begin{proof}
  This is the declared model definition of Acoustic Frequency.
\end{proof}
\begin{definition}[Acoustic Speed]
  \label{def:physics:phyx-mini-0186:phyxminiproblems-problemphyxmini0186-acousticspeed}
  \lean{PhyXMiniProblems.ProblemPhyXMini0186.AcousticSpeed}
  A physical sound speed, carrying the length-per-time dimension
\end{definition}
\begin{proof}
  This is the declared model definition of Acoustic Speed.
\end{proof}
\begin{definition}[meters Value]
  \label{def:physics:phyx-mini-0186:phyxminiproblems-problemphyxmini0186-metersvalue}
  \lean{PhyXMiniProblems.ProblemPhyXMini0186.metersValue}
  \uses{def:physics:phyx-mini-0186:phyxminiproblems-problemphyxmini0186-acousticlength}
  Read a dimensionful length in metres
\end{definition}
\begin{proof}
  This is the declared model definition of meters Value.
\end{proof}
\begin{definition}[hertz Value]
  \label{def:physics:phyx-mini-0186:phyxminiproblems-problemphyxmini0186-hertzvalue}
  \lean{PhyXMiniProblems.ProblemPhyXMini0186.hertzValue}
  \uses{def:physics:phyx-mini-0186:phyxminiproblems-problemphyxmini0186-acousticfrequency}
  Read a dimensionful acoustic frequency in hertz
\end{definition}
\begin{proof}
  This is the declared model definition of hertz Value.
\end{proof}
\begin{definition}[meters Per Second Value]
  \label{def:physics:phyx-mini-0186:phyxminiproblems-problemphyxmini0186-meterspersecondvalue}
  \lean{PhyXMiniProblems.ProblemPhyXMini0186.metersPerSecondValue}
  \uses{def:physics:phyx-mini-0186:phyxminiproblems-problemphyxmini0186-acousticspeed}
  Read a dimensionful sound speed in metres per second
\end{definition}
\begin{proof}
  This is the declared model definition of meters Per Second Value.
\end{proof}
\begin{definition}[length Of Meters]
  \label{def:physics:phyx-mini-0186:phyxminiproblems-problemphyxmini0186-lengthofmeters}
  \lean{PhyXMiniProblems.ProblemPhyXMini0186.lengthOfMeters}
  \uses{def:physics:phyx-mini-0186:phyxminiproblems-problemphyxmini0186-acousticlength}
  Construct a physical length from its metre readout
\end{definition}
\begin{proof}
  This is the declared model definition of length Of Meters.
\end{proof}
\begin{definition}[Speaker Label]
  \label{def:physics:phyx-mini-0186:phyxminiproblems-problemphyxmini0186-speakerlabel}
  \lean{PhyXMiniProblems.ProblemPhyXMini0186.SpeakerLabel}
  The loudspeaker numbers printed in the primary figure
\end{definition}
\begin{proof}
  This is the declared model definition of Speaker Label.
\end{proof}
\begin{definition}[Three Speaker Interference Setup]
  \label{def:physics:phyx-mini-0186:phyxminiproblems-problemphyxmini0186-threespeakerinterferencesetup}
  \lean{PhyXMiniProblems.ProblemPhyXMini0186.ThreeSpeakerInterferenceSetup}
  \uses{def:physics:phyx-mini-0186:phyxminiproblems-problemphyxmini0186-acousticlength, def:physics:phyx-mini-0186:phyxminiproblems-problemphyxmini0186-acousticfrequency, def:physics:phyx-mini-0186:phyxminiproblems-problemphyxmini0186-acousticspeed, def:physics:phyx-mini-0186:phyxminiproblems-problemphyxmini0186-speakerlabel}
  The physical objects and named quantities in the three-speaker experiment. The coordinates of speakerPosition and listenerPosition are metre readouts in the axes of the primary figure: positive x points toward the listener and positive y points toward speaker 1. The common individual wave amplitude at the listener is the quantity denoted by a in the problem
\end{definition}
\begin{proof}
  This is the declared model definition of Three Speaker Interference Setup.
\end{proof}
\begin{definition}[Matches Problem Readouts]
  \label{def:physics:phyx-mini-0186:phyxminiproblems-problemphyxmini0186-matchesproblemreadouts}
  \lean{PhyXMiniProblems.ProblemPhyXMini0186.MatchesProblemReadouts}
  \uses{def:physics:phyx-mini-0186:phyxminiproblems-problemphyxmini0186-metersvalue, def:physics:phyx-mini-0186:phyxminiproblems-problemphyxmini0186-hertzvalue, def:physics:phyx-mini-0186:phyxminiproblems-problemphyxmini0186-meterspersecondvalue, def:physics:phyx-mini-0186:phyxminiproblems-problemphyxmini0186-threespeakerinterferencesetup}
  The scalar data printed in the problem and figure: adjacent speakers are 3 m apart, the listener is 4 m in front of speaker 2, the common tone is 170 Hz, and sound travels at 340 m/s. This predicate contains neither a leftward displacement nor a maximum-interference answer
\end{definition}
\begin{proof}
  This is the declared model definition of Matches Problem Readouts.
\end{proof}
\begin{definition}[Has Depicted Layout]
  \label{def:physics:phyx-mini-0186:phyxminiproblems-problemphyxmini0186-hasdepictedlayout}
  \lean{PhyXMiniProblems.ProblemPhyXMini0186.HasDepictedLayout}
  \uses{def:physics:phyx-mini-0186:phyxminiproblems-problemphyxmini0186-metersvalue, def:physics:phyx-mini-0186:phyxminiproblems-problemphyxmini0186-threespeakerinterferencesetup}
  The coordinate layout read from the primary image. Speaker 2 is the origin; speakers 1 and 3 lie respectively 3 m above and below it, and the listener lies on the positive horizontal axis
\end{definition}
\begin{proof}
  This is the declared model definition of Has Depicted Layout.
\end{proof}
\begin{definition}[Emits Coherently With Common Amplitude]
  \label{def:physics:phyx-mini-0186:phyxminiproblems-problemphyxmini0186-emitscoherentlywithcommonamplitude}
  \lean{PhyXMiniProblems.ProblemPhyXMini0186.EmitsCoherentlyWithCommonAmplitude}
  \uses{def:physics:phyx-mini-0186:phyxminiproblems-problemphyxmini0186-threespeakerinterferencesetup}
  The identical-source assumptions used by the interference model. All three sources emit with one initial phase, and the wave arriving from each source has the named common amplitude a, as stated in the problem. Their frequency is already common because it is represented by the single field toneFrequency
\end{definition}
\begin{proof}
  This is the declared model definition of Emits Coherently With Common Amplitude.
\end{proof}
\begin{definition}[Has Physical Parameters]
  \label{def:physics:phyx-mini-0186:phyxminiproblems-problemphyxmini0186-hasphysicalparameters}
  \lean{PhyXMiniProblems.ProblemPhyXMini0186.HasPhysicalParameters}
  \uses{def:physics:phyx-mini-0186:phyxminiproblems-problemphyxmini0186-metersvalue, def:physics:phyx-mini-0186:phyxminiproblems-problemphyxmini0186-hertzvalue, def:physics:phyx-mini-0186:phyxminiproblems-problemphyxmini0186-meterspersecondvalue, def:physics:phyx-mini-0186:phyxminiproblems-problemphyxmini0186-threespeakerinterferencesetup}
  Positivity conditions for the physical wavelength and supplied data
\end{definition}
\begin{proof}
  This is the declared model definition of Has Physical Parameters.
\end{proof}
\begin{definition}[Satisfies Sound Wave Law]
  \label{def:physics:phyx-mini-0186:phyxminiproblems-problemphyxmini0186-satisfiessoundwavelaw}
  \lean{PhyXMiniProblems.ProblemPhyXMini0186.SatisfiesSoundWaveLaw}
  \uses{def:physics:phyx-mini-0186:phyxminiproblems-problemphyxmini0186-metersvalue, def:physics:phyx-mini-0186:phyxminiproblems-problemphyxmini0186-hertzvalue, def:physics:phyx-mini-0186:phyxminiproblems-problemphyxmini0186-meterspersecondvalue, def:physics:phyx-mini-0186:phyxminiproblems-problemphyxmini0186-threespeakerinterferencesetup}
  The acoustic propagation law c = λ f. It is a governing law relating speed, wavelength, and frequency, not a statement about the requested motion of speaker 2
\end{definition}
\begin{proof}
  This is the declared model definition of Satisfies Sound Wave Law.
\end{proof}
\begin{definition}[moved Middle Speaker Position]
  \label{def:physics:phyx-mini-0186:phyxminiproblems-problemphyxmini0186-movedmiddlespeakerposition}
  \lean{PhyXMiniProblems.ProblemPhyXMini0186.movedMiddleSpeakerPosition}
  \uses{def:physics:phyx-mini-0186:phyxminiproblems-problemphyxmini0186-acousticlength, def:physics:phyx-mini-0186:phyxminiproblems-problemphyxmini0186-metersvalue, def:physics:phyx-mini-0186:phyxminiproblems-problemphyxmini0186-threespeakerinterferencesetup}
  The position obtained by moving speaker 2 left by displacement. Since the listener lies on the positive horizontal axis, a leftward movement subtracts the positive metre readout from speaker 2's x coordinate and leaves its vertical coordinate fixed
\end{definition}
\begin{proof}
  This is the declared model definition of moved Middle Speaker Position.
\end{proof}
\begin{definition}[source Position After Move]
  \label{def:physics:phyx-mini-0186:phyxminiproblems-problemphyxmini0186-sourcepositionaftermove}
  \lean{PhyXMiniProblems.ProblemPhyXMini0186.sourcePositionAfterMove}
  \uses{def:physics:phyx-mini-0186:phyxminiproblems-problemphyxmini0186-acousticlength, def:physics:phyx-mini-0186:phyxminiproblems-problemphyxmini0186-speakerlabel, def:physics:phyx-mini-0186:phyxminiproblems-problemphyxmini0186-threespeakerinterferencesetup, def:physics:phyx-mini-0186:phyxminiproblems-problemphyxmini0186-movedmiddlespeakerposition}
  Source positions after moving only speaker 2 to the left
\end{definition}
\begin{proof}
  This is the declared model definition of source Position After Move.
\end{proof}
\begin{definition}[path Length Meters After Move]
  \label{def:physics:phyx-mini-0186:phyxminiproblems-problemphyxmini0186-pathlengthmetersaftermove}
  \lean{PhyXMiniProblems.ProblemPhyXMini0186.pathLengthMetersAfterMove}
  \uses{def:physics:phyx-mini-0186:phyxminiproblems-problemphyxmini0186-acousticlength, def:physics:phyx-mini-0186:phyxminiproblems-problemphyxmini0186-speakerlabel, def:physics:phyx-mini-0186:phyxminiproblems-problemphyxmini0186-threespeakerinterferencesetup, def:physics:phyx-mini-0186:phyxminiproblems-problemphyxmini0186-sourcepositionaftermove}
  The propagation distance from one speaker to the listener after the move, read in metres through the metric on Physlib's planar Space 2
\end{definition}
\begin{proof}
  This is the declared model definition of path Length Meters After Move.
\end{proof}
\begin{definition}[Produces Maximum Amplitude At Listener]
  \label{def:physics:phyx-mini-0186:phyxminiproblems-problemphyxmini0186-producesmaximumamplitudeatlistener}
  \lean{PhyXMiniProblems.ProblemPhyXMini0186.ProducesMaximumAmplitudeAtListener}
  \uses{def:physics:phyx-mini-0186:phyxminiproblems-problemphyxmini0186-acousticlength, def:physics:phyx-mini-0186:phyxminiproblems-problemphyxmini0186-metersvalue, def:physics:phyx-mini-0186:phyxminiproblems-problemphyxmini0186-speakerlabel, def:physics:phyx-mini-0186:phyxminiproblems-problemphyxmini0186-threespeakerinterferencesetup, def:physics:phyx-mini-0186:phyxminiproblems-problemphyxmini0186-pathlengthmetersaftermove}
  For coherent monochromatic sources of equal individual amplitude, their sum has maximum amplitude precisely when every pairwise path difference is an integer multiple of the common wavelength. This predicate records that standard constructive-interference condition for a proposed displacement; it does not specify which displacement satisfies it
\end{definition}
\begin{proof}
  This is the declared model definition of Produces Maximum Amplitude At Listener.
\end{proof}
\begin{definition}[Is Least Positive Maximum Displacement]
  \label{def:physics:phyx-mini-0186:phyxminiproblems-problemphyxmini0186-isleastpositivemaximumdisplacement}
  \lean{PhyXMiniProblems.ProblemPhyXMini0186.IsLeastPositiveMaximumDisplacement}
  \uses{def:physics:phyx-mini-0186:phyxminiproblems-problemphyxmini0186-acousticlength, def:physics:phyx-mini-0186:phyxminiproblems-problemphyxmini0186-metersvalue, def:physics:phyx-mini-0186:phyxminiproblems-problemphyxmini0186-threespeakerinterferencesetup, def:physics:phyx-mini-0186:phyxminiproblems-problemphyxmini0186-producesmaximumamplitudeatlistener}
  A displacement is the requested move when it is positive, produces a constructive maximum, and is no larger than any other positive displacement that produces a maximum. This makes the textbook question's intended "first maximum" explicit and rules out later maxima one or more wavelengths farther to the left
\end{definition}
\begin{proof}
  This is the declared model definition of Is Least Positive Maximum Displacement.
\end{proof}
\begin{lemma}[wavelength in meters eq two]
  \label{lem:physics:phyx-mini-0186:phyxminiproblems-problemphyxmini0186-wavelength-in-meters-eq-two}
  \lean{PhyXMiniProblems.ProblemPhyXMini0186.wavelength_in_meters_eq_two}
  \uses{def:physics:phyx-mini-0186:phyxminiproblems-problemphyxmini0186-metersvalue, def:physics:phyx-mini-0186:phyxminiproblems-problemphyxmini0186-threespeakerinterferencesetup, def:physics:phyx-mini-0186:phyxminiproblems-problemphyxmini0186-matchesproblemreadouts, def:physics:phyx-mini-0186:phyxminiproblems-problemphyxmini0186-hasphysicalparameters, def:physics:phyx-mini-0186:phyxminiproblems-problemphyxmini0186-satisfiessoundwavelaw}
  The numerical sound data and c = λf imply a wavelength of 2 m
\end{lemma}
\begin{proof}
  The conclusion follows from the governing relations and figure readouts named in the statement of wavelength in meters eq two.
\end{proof}
\begin{lemma}[initial path lengths are five four five]
  \label{lem:physics:phyx-mini-0186:phyxminiproblems-problemphyxmini0186-initial-path-lengths-are-five-four-five}
  \lean{PhyXMiniProblems.ProblemPhyXMini0186.initial_path_lengths_are_five_four_five}
  \uses{def:physics:phyx-mini-0186:phyxminiproblems-problemphyxmini0186-lengthofmeters, def:physics:phyx-mini-0186:phyxminiproblems-problemphyxmini0186-threespeakerinterferencesetup, def:physics:phyx-mini-0186:phyxminiproblems-problemphyxmini0186-matchesproblemreadouts, def:physics:phyx-mini-0186:phyxminiproblems-problemphyxmini0186-hasdepictedlayout, def:physics:phyx-mini-0186:phyxminiproblems-problemphyxmini0186-pathlengthmetersaftermove}
  Before speaker 2 moves, the 3-4-5 triangles in the figure give path lengths 5 m, 4 m, and 5 m from speakers 1, 2, and 3 respectively
\end{lemma}
\begin{proof}
  The conclusion follows from the governing relations and figure readouts named in the statement of initial path lengths are five four five.
\end{proof}
\begin{lemma}[one meter move equalizes path lengths]
  \label{lem:physics:phyx-mini-0186:phyxminiproblems-problemphyxmini0186-one-meter-move-equalizes-path-lengths}
  \lean{PhyXMiniProblems.ProblemPhyXMini0186.one_meter_move_equalizes_path_lengths}
  \uses{def:physics:phyx-mini-0186:phyxminiproblems-problemphyxmini0186-lengthofmeters, def:physics:phyx-mini-0186:phyxminiproblems-problemphyxmini0186-threespeakerinterferencesetup, def:physics:phyx-mini-0186:phyxminiproblems-problemphyxmini0186-matchesproblemreadouts, def:physics:phyx-mini-0186:phyxminiproblems-problemphyxmini0186-hasdepictedlayout, def:physics:phyx-mini-0186:phyxminiproblems-problemphyxmini0186-pathlengthmetersaftermove}
  Moving speaker 2 one metre left puts it five metres from the listener, equal to the two unchanged outer-speaker path lengths
\end{lemma}
\begin{proof}
  The conclusion follows from the governing relations and figure readouts named in the statement of one meter move equalizes path lengths.
\end{proof}
\begin{definition}[Answer Choice]
  \label{def:physics:phyx-mini-0186:phyxminiproblems-problemphyxmini0186-answerchoice}
  \lean{PhyXMiniProblems.ProblemPhyXMini0186.AnswerChoice}
  Labels of the four answer choices printed with the problem
\end{definition}
\begin{proof}
  This is the declared model definition of Answer Choice.
\end{proof}
\begin{definition}[meters]
  \label{def:physics:phyx-mini-0186:phyxminiproblems-problemphyxmini0186-answerchoice-meters}
  \lean{PhyXMiniProblems.ProblemPhyXMini0186.AnswerChoice.meters}
  \uses{def:physics:phyx-mini-0186:phyxminiproblems-problemphyxmini0186-answerchoice}
  The displacement in metres printed beside each answer label
\end{definition}
\begin{proof}
  This is the declared model definition of meters.
\end{proof}
\begin{definition}[Is Correct Answer]
  \label{def:physics:phyx-mini-0186:phyxminiproblems-problemphyxmini0186-iscorrectanswer}
  \lean{PhyXMiniProblems.ProblemPhyXMini0186.IsCorrectAnswer}
  \uses{def:physics:phyx-mini-0186:phyxminiproblems-problemphyxmini0186-acousticlength, def:physics:phyx-mini-0186:phyxminiproblems-problemphyxmini0186-metersvalue, def:physics:phyx-mini-0186:phyxminiproblems-problemphyxmini0186-threespeakerinterferencesetup, def:physics:phyx-mini-0186:phyxminiproblems-problemphyxmini0186-isleastpositivemaximumdisplacement, def:physics:phyx-mini-0186:phyxminiproblems-problemphyxmini0186-answerchoice}
  A displayed choice is correct when it gives the least positive maximum
\end{definition}
\begin{proof}
  This is the declared model definition of Is Correct Answer.
\end{proof}
\begin{definition}[recorded Answer Choice]
  \label{def:physics:phyx-mini-0186:phyxminiproblems-problemphyxmini0186-recordedanswerchoice}
  \lean{PhyXMiniProblems.ProblemPhyXMini0186.recordedAnswerChoice}
  \uses{def:physics:phyx-mini-0186:phyxminiproblems-problemphyxmini0186-answerchoice}
  The answer label recorded in the supplied dataset metadata
\end{definition}
\begin{proof}
  This is the declared model definition of recorded Answer Choice.
\end{proof}
% --- Archon declaration topology end ---
% --- Archon physics formalization source end ---
